\section{Source reconciliation, verification, and further research}
\label{sec:audit}

\subsection{Which earlier questions are settled}

The following table records the precise scope of the continuation.
The source labels refer to the incoming reports at the pinned
snapshot, not to a moving version of the repository.

\begin{longtable}{@{}>{\raggedright\arraybackslash}p{0.27\textwidth}>{\raggedright\arraybackslash}p{0.65\textwidth}@{}}
\toprule
\textbf{Earlier target} & \textbf{Result and remaining scope}\\
\midrule
\endhead
Vertical lattice, Questions 1 and 6 \cite{RepoVertical}
&
Theorems~\ref{thm:minimal-contact} and \ref{thm:filter-all}
give every minimal contact and the exact periodic mean criterion.
All derivative orders are covered.\\[5pt]
Vertical lattice, Question 7 \cite{RepoVertical}
&
Theorem~\ref{thm:modular} gives exact transport for horizontal
shift \(a=1\), near every nonpositive even order and near \(s=1\),
with every derivative correction. A general arbitrary-shift
transport remains open.\\[5pt]
Unequal products beyond two factors \cite{RepoBernoulli,RepoMixed}
&
Theorem~\ref{mtp:generator} handles every finite collection of
positive even tail orders and every regular centered moment.
The separate conjecture about optimal unequal bilinear
prefixes remains open.\\[5pt]
Independent nonperiodic Stieltjes shifts, R8 \cite{RepoTailKernels}
&
Theorem~\ref{gs:completed} supplies an explicit jointly
holomorphic convergent generator, including all argument
derivatives. A finite reduction in a prescribed classical
constant alphabet is a stronger question.\\[5pt]
Higher spectral collisions and independent orders
\cite{RepoBernoulli}
&
The face-moment formula gives all polynomial collision
coefficients. The one-displacement Gamma sector is evaluated
by depth-two polylogarithms. The joint polar part and its
finite local correction are explicit.\\[5pt]
Alternating nonpositive-order square law, R7
\cite{RepoTailKernels}
&
Already resolved in the separate Bilinear Tails archive
\cite{RepoBilinear}; update its status and cite that
proof. It is not counted as a new result here.\\[5pt]
Gaussian \(S_6\) and revised \(S_8\)
\cite{RepoManuscript}
&
Remain conjectural. The present results do not supply the
missing period relations.\\
\bottomrule
\end{longtable}

\subsection{A concrete correction to the research ledger}

In the Tail Kernels archive, R7 asks for the alternating square
continuation at all nonpositive orders. At the same snapshot,
\texttt{sections/04\_cutoff\_jets.tex} in the Bilinear Tails
archive already proves the required normally convergent
continuation, full contact formula, and translated-cutoff rule
\cite{RepoTailKernels,RepoBilinear}. Its theorem labels
\texttt{cut:normal}, \texttt{cut:contacts}, and
\texttt{cut:translation} locate the exact claims.
The order-zero value \(1/4\), despite the vanishing
pointwise summands in the relevant representation, is explicitly
reconciled there. The Gamma specializations have labels
\texttt{cut:Gamma} and \texttt{cut:Gammasquare}.

The proposed replacement ledger entry is:
\begin{quote}
\textbf{R7: resolved.} Cite the all-order cutoff continuation in
\emph{Bilinear Alternating Tails and Harmonic Determinants},
Section 4. Preserve the specified cutoff coordinate and its
order-zero contact. Remaining extensions should be stated
separately from the already established nonpositive-order law.
\end{quote}
This corrects the status shared between two reports. It does not
identify a false theorem in the resolving report.

\subsection{Three safeguards for subsequent integration}

\paragraph{Finite parts must retain their coordinate.}
For the new lattice formula, the constant left after deleting
positive powers of \(\log K\) is
\({\cal F}_M^{(d)}(-N)+C_{N,d}(\log h)\).
Deleting positive powers of the physical-radius logarithm
\(\log R\), \(R=hK\), instead gives
\({\cal F}_M^{(d)}(-N)+C_{N,d}(0)\).
The analytic continuation value is obtained only after the entire
contact polynomial has been removed. These are three specified
operations; their constant differences should not be suppressed
in a shared definition of \emph{regularized sum}.

\paragraph{Periodic cancellation does not imply absolute convergence.}
A nonzero mean-zero periodic multiplier removes the critical
contact at every fixed derivative order. The resulting two
half-lattice series converge by summation by parts. They are in
general conditional, since the terms still have critical
power-log magnitudes. An additional horizontal difference
improves the decay and gives absolute convergence.
Theorem~\ref{thm:filter-all} and
\eqref{eq:relative-lattice} keep these claims distinct.

\paragraph{A common-ray limit is insufficient for independent orders.}
The predecessor's common-ray spectral theorem is correctly
scoped. An unqualified extension to independent orders would
be false. For instance, with \(W(x)=x\), \(a=1\), and masses
\(0,1\), the limits along \((\alpha_1,\alpha_2)=(t,t)\) and
\((t,2t)\) are \(-1/3\) and \(-5/12\).
Even equality of all nontangential straight-ray limits would
not suffice: the example \(W(x)=x^3+x/2\) has polar part
\[
 \frac{\alpha_2^2}{4(\alpha_1+\alpha_2)}.
\]
It tends to zero on every fixed ray with nonzero total direction,
but along \(\alpha_2=t\),
\(\alpha_1=-t+t^3\) it equals \(1/(4t)\).
The joint polar decomposition, rather than a ray calculation,
is the appropriate hypothesis for mixed differentiation.

\subsection{Reproducible verification}

The companion files separate exact algebra from numerical
diagnostics. All arithmetic determinants and finite polynomial
identities are computed over rational or symbolic exact rings.
The floating-point checks compare independently organized
expressions, including direct quadratures, finite cutoff sums,
and Mellin contours.

\begin{longtable}{@{}>{\raggedright\arraybackslash}p{0.25\textwidth}>{\raggedright\arraybackslash}p{0.67\textwidth}@{}}
\toprule
\textbf{Component} & \textbf{Verification supplied}\\
\midrule
\endhead
Lattice contacts
&
60 exact coefficient and Pochhammer checks, two mutation
controls, and 27 finite-lattice diagnostics at 55 digits.
The next two Euler--Maclaurin coefficients are used to check
the remainder when its leading polynomial nearly cancels.\\[4pt]
Modular transport
&
Nine exact kernel-polynomial checks, four direct
Lambert-versus-contour comparisons at 42 digits,
two classical boundary checks, and an independent
self-dual contour evaluation.\\[4pt]
Mixed tails
&
306 exact moments for all multisets of one through four
orders from \(\{2,4,6\}\), through the tested moment range,
against finite summation by parts; 825 boundary-parity
checks; exact six- and eight-coordinate determinants;
a finite telescoping certificate; and 90-digit checks
of the short identities.\\[4pt]
Independent shifts
&
16 exact cancellation checks and 42 numerical checks
at 65 digits, including complex shifts, both spectral
axes, \(u+v=0\), prefix invariance, and the explicit
\(\gamma_1\times\psi\) coefficient.
The largest recorded absolute residual is approximately
\(1.02\times10^{-56}\).\\[4pt]
Spectral faces
&
34 exact rational checks and 11 numerical comparisons
at 48-digit working precision. The original
two-dimensional Gamma face is checked against the
polylogarithmic lattice expression with increasing
quadrature orders.\\
\bottomrule
\end{longtable}

These calculations are regression diagnostics for the stated
proofs. They do not establish global arithmetic independence,
replace a normal-convergence argument, or certify decimal
enclosures. All domains, branch conventions, and limiting
operations used in the identities are established in the text.
The verification programs and their recorded outputs are
included so that future integration can detect sign, endpoint,
normalization, or coefficient regressions.

\subsection{Further research questions}
\label{sec:questions}

The following questions are ordered by their connection to the
new mechanisms and by the prospect of obtaining further exact
identities. Each asks for a structural statement or evaluation.

\begin{enumerate}[label=\textbf{Q\arabic*.},leftmargin=*]
\item \textbf{Optimal prefixes for arbitrary unequal tails.}
The mixed recurrence proves a finite span, and the two worked
products have complete rational prefix certificates.
Determine the minimal formal coordinate dimension for a
general even-order multiset \(I\), after a declared system of
MZV reductions, and determine when its first admissible
moments already span that space. The existing unequal
bilinear optimal-prefix conjecture is a particularly precise
first target. A proof should establish an explicit nonzero
determinant, rather than infer rank from evaluated decimals.

\item \textbf{Direct telescoping certificates for elegant tail identities.}
The short cubic identity admits a finite telescoping
certificate in products of Hurwitz tails. Characterize the
centered polynomial combinations for which such a certificate
exists. Can the polarized recurrence be compiled into a
finite certificate with a prescribed maximal tail order?
This could systematically produce short ordinary identities
whose cancellations hide a much larger moment system.

\item \textbf{Minimal functional depth of mixed generators.}
The bound \(d-1\) is constructive but need not be sharp.
Classify the order multisets for which the anchored generator
reduces to classical polylogarithms, or to depth two
independently of \(d\). The classification must distinguish
functional identities from numerical MZV reductions at the
endpoint \(t=0\).

\item \textbf{More than one displaced spectral mass.}
For \(k-1\) equal masses and one displaced mass, depth two
suffices in every dimension. Determine the exact function
class for two displaced masses and, more generally, for a
fixed number of distinct mass clusters. Is depth controlled
by the number of clusters rather than the number of factors?
A useful result would give a finite multiple-polylogarithm
formula with explicit branches and ordinary convergence,
not merely rename the remaining simplex integral.

\item \textbf{Complete logarithmic collision coefficients.}
The rational normalized face formula produces every
logarithmically weighted polynomial moment by differentiation.
Find compact closed forms for the first genuinely nonsquarefree
mixed order derivatives, and organize their coefficients by
harmonic star sums and zeta weight. Determine which Gamma
normalization coefficients cancel before any MZV reduction.

\item \textbf{Independent-order continuation beyond quadratic spectra.}
For quadratic factors, the odd powers of \(W\) give the exact
polar obstruction through a single total-order divisor.
Develop the corresponding finite polar-part classification
for higher-degree polynomial spectra and shifted linear
factors. Identify the natural local subtraction that is
compatible with multiplication, parameter differentiation,
and periodic mean cancellation.

\item \textbf{Modular transport with arbitrary horizontal shift.}
Extend the exact order-dependent contour identity from
\(a=1\) to general \(a\). Rational shifts should first be
organized into finite systems with additive characters and
Hurwitz factors. For irrational shifts, determine the
additional integral transform explicitly. The fixed-cutoff
contact polynomial should remain visible throughout the
transport.

\item \textbf{Arithmetic reductions of the logarithmic Hurwitz family.}
For which rational pairs \((c,z)\) does
\[
 \mathcal L_p(c,z)=
 \sum_{n\ge0}\frac{\log(1+z/(n+c))}{(n+c)^p}
\]
reduce to a finite specified alphabet of polylogarithms,
Gamma values, and zeta derivatives?
The generating and differential relations in the
independent-shift section impose strong compatibility
conditions. Generic rational arguments should not be
assumed to have such a reduction.

\item \textbf{Unequal-rate Stieltjes boundary collisions.}
Set \(a=\lambda\epsilon\), \(b=\mu\epsilon\) with fixed
\(\lambda,\mu>0\), and let \(\epsilon\downarrow0\).
Compute the complete singular power-log primitive and
the remaining constant for arbitrary Stieltjes indices
and argument derivatives. Determine the exact path
correction relative to the common-shift finite part.
The diagonal unfolding and explicit finite-prefix
derivatives provide a starting point.

\item \textbf{Three independently shifted Stieltjes factors.}
Extend the evaluated two-factor germ to three factors
with independent shifts and orders. Partitioning the
summation lattice into minimum-index regions suggests
finite Lauricella-type prefixes. The central problem is
to remove all intersecting spectral poles jointly while
retaining geometric convergence after every derivative.

\item \textbf{The unresolved Gaussian period identities.}
The manuscript's \(S_6\) and revised \(S_8\) candidates
need exact period relations. The existing high-precision
residuals and bounded relation-space calculations do
not supply them. Seek a new functional relation or
contour deformation that escapes the already tested
relation spaces, and exhibit its exact coefficient
certificate before numerical evaluation. The new
mixed-tail rank results concern different spaces and
should not be used as evidence of a Gaussian proof.

\item \textbf{A formal proof layer for the finite core.}
The nilpotent polarized Bernoulli identity, the exact
boundary-parity cancellation, the contact-polynomial
coefficient identity, and the rational face compiler
are natural Lean targets. Separate their finite
algebraic statements from Euler--Maclaurin continuation
and normal-convergence lemmas. This would make later
analytic generalizations depend on a small reusable
verified algebraic core, without claiming that the
current Python checks constitute formal proofs.
\end{enumerate}

